\documentclass[11pt,a4paper]{article}
\usepackage[utf8]{inputenc}
\usepackage[T1]{fontenc}
\usepackage[french]{babel}
\usepackage{geometry}
\usepackage{xcolor}
\usepackage{hyperref}
\usepackage{titlesec}

\geometry{margin=1in}
\definecolor{titleblue}{RGB}{0,51,102}

\hypersetup{
    colorlinks=true,
    linkcolor=titleblue,
    urlcolor=titleblue,
    pdftitle={Lettre de Motivation - Ismail AISSAOUI},
}

\begin{document}

\noindent \textbf{Ismail AISSAOUI} \\
Ingénieur d'État en Intelligence Artificielle (ESI-SBA) \\
Email: ismail.aissaoui.pro@gmail.com \\
Téléphone: +213 660 70 77 96 \\
LinkedIn: https://ismail-aissaoui.vercel.app

\bigskip

\begin{flushright}
    Au Comité de Sélection du Doctorat \\
    \textbf{Programme EDT - INRAE / Inria} \\
    À l'attention de : Dr. Lorenzo Sala et Prof. Julien Deantoni \\
    \medskip
    1er mai 2026
\end{flushright}

\bigskip

\noindent \textbf{Objet : Candidature au doctorat : Jumeaux numériques hybrides compositionnels pour les écosystèmes microbiens fermentés (Rang 7 - PC1)}

\bigskip

Madame, Monsieur,

Ingénieur d'État en Intelligence Artificielle diplômé de l'ESI-SBA, je vous adresse avec enthousiasme ma candidature pour le poste de doctorat intitulé « Jumeaux numériques hybrides compositionnels pour les écosystèmes microbiens fermentés » au sein du programme national Engineering Digital Twin (EDT). Cette candidature est motivée par la forte adéquation de mon profil avec votre projet interdisciplinaire.

Actuellement en fin de cycle ingénieur chez Sonatrach, je travaille sur la **hybridation de modèles mécanistes et de substituts neuronaux** pour la surveillance de systèmes industriels. J'ai acquis une solide expérience dans la conception de fonctions de perte **PINN (Physics-Informed Neural Network)** garantissant la cohérence physique des modèles. Je suis particulièrement motivé par l'idée de transposer cette expertise au monde complexe de la fermentation microbienne et d'investiguer des opérateurs formels pour la composition de modèles et le mélange multi-fidélité.

Votre projet à l'INRAE et à l'Inria/UniCA m'attire vivement car il propose de passer de simulations ad-hoc à un cadre de modélisation hybride rigoureux pour la biologie. Je suis convaincu que ma maîtrise de PyTorch et ma compréhension des systèmes dynamiques (EDO) me permettront de contribuer de manière significative à vos objectifs de développement de jumeaux numériques fiables et scalables.

Je suis un chercheur autonome et passionné par les défis interdisciplinaires, prêt à m'investir dans la compréhension des mécanismes biologiques pour enrichir la **plateforme Artemis**.

Je vous remercie de l'attention que vous porterez à ma candidature et reste à votre entière disposition pour un entretien.

Dans l'attente de votre réponse, je vous prie d'agréer, Madame, Monsieur, l'expression de mes salutations distinguées.

\bigskip

\noindent Ismail AISSAOUI

\end{document}
